\section{Discussion}
\label{sec:discussion}

Predicting the assembly, succession, and long-term stability of microbial communities has traditionally been hindered by a dichotomy between overly complex, phenomenological interaction models (e.g., high-dimensional Generalized Lotka-Volterra equations) and overly simplistic discrete heuristics that ignore intra-passage growth kinetics. In this study, we demonstrated that a \textbf{Physical Consumer-Resource Neural ODE} provides a parsimonious, generalizable, and mathematically rigorous foundation for serial dilution ecology across vastly different biological systems.

\subsection{Overcoming the Competitive Exclusion Principle}
Classical ecological theory predicts that in a well-mixed environment with $K$ limiting resources, at most $K$ species can stably coexist at equilibrium (the Competitive Exclusion Principle) \citep{macarthur1970species, tilman1982resource}. Yet empirical communities regularly harbor dozens of coexisting taxa even in simple media \citep{goldford2018emergent, hu2022emergent}. In our framework, multi-species coexistence is sustained by two physical mechanisms:
\begin{enumerate}
    \item \textbf{Quadratic Self-Limitation ($-\alpha_i N_i^2$):} By penalizing high intra-strain densities through space saturation, phage susceptibility, or autotoxicity, self-limitation prevents the fastest-growing taxon from monopolizing resources, thereby opening temporal and spatial niches for slower-growing competitors.
    \item \textbf{Non-Equilibrium Periodic Forcing:} The periodic dilution transfer ($D$) resets communities to low biomass before steady-state competitive exclusion can be finalized, allowing species with complementary growth rates and yields to persist indefinitely.
\end{enumerate}

\subsection{The Necessity of Physical Biomass Continuity}
A key technical finding of this study is that normalizing abundance vectors to the unit simplex at transfer boundaries severely distorts ecological simulations. Because quadratic penalties scale with the square of biomass ($N_i^2$), artificially compressing biomass into relative fractions ($x_i \approx 0.02$) attenuates quadratic regulation by up to four orders of magnitude ($x_i^2 \approx 0.0004$). Restoring authentic unnormalized biomass propagation across serial transfers restored physiological self-limitation, eliminating artificial extinction cascades and directly boosting predictive accuracy on empirical datasets by $+4.4\%$ (Pollard) and $+9.4\%$ (German).

\subsection{Minimalist Latent Niches \texorpdfstring{($K \ll N$)}{(K << N)} and Scalability}
Across all four datasets, low-dimensional latent resource spaces ($K=2$ to $K=12$) were sufficient to achieve high predictive accuracy ($R^2 = 0.73\text{--}0.97$). Rather than explicitly enumerating dozens of chemical metabolites, the model projects metabolic interactions onto orthogonal effective niches, capturing primary substrate utilization and cross-feeding networks with minimal parameter counts. This compactness prevents overfitting, provides interpretable environmental representations, and enables zero-shot generalizations to unseen community mixtures and nutrient regimes.

\subsection{From Physical Resources to Emergent, Environment-Dependent Interaction Networks}
A long-standing debate in microbial ecology concerns whether community dynamics should be modeled via direct species-species interactions (e.g., Generalized Lotka-Volterra) or through explicit environmental resource mediation. Our analytical derivation (Section~\ref{sec:interaction_matrix}) reconciles these two paradigms by demonstrating that the effective interaction matrix $\mathbf{A}^{(c)}$ is an emergent, context-dependent operator rather than a collection of static parameters.

Crucially, our findings show that pairwise interactions are not immutable constants:
\begin{itemize}
    \item \textbf{Environmental Conditioning:} Because $\mathbf{A}^{(c)}$ depends directly on the external resource vector $\mathbf{R}_{0, c}$, changing nutrient availability (such as glucose and urea in Gore et al., or spatial carriers in Pollard et al.) physically re-tunes the entire interaction network. When resources are elevated, interaction strengths scale up, driving high-richness communities across dynamical bifurcations into chaos.
    \item \textbf{Sign Reversals:} The net interaction between two species can flip between mutualistic cross-feeding and antagonistic competition depending on whether secondary metabolites or primary substrates are limiting in the medium. This context-dependence explains why pairwise co-culture interactions measured in monoculture broths frequently fail to predict community assembly in complex backgrounds \citep{venturelli2018deciphering, friedman2017community}.
    \item \textbf{Resolution of the May-Wigner Paradox:} Robert May's classic stability criterion states that random, fully connected interaction networks become mathematically unstable when complexity $N \sqrt{C} \sigma > 1$ \citep{may1972will}. In our framework, the effective interaction matrix is strictly constrained by a low-dimensional metabolic manifold ($\text{rank} \le K \ll N$). This intrinsic low-rank constraint breaks random all-to-all feedback loops, providing a clear biophysical mechanism for how natural microbiomes maintain high taxonomic richness without succumbing to dynamical collapse.
\end{itemize}
